\documentclass[11pt]{article}

\usepackage[utf8]{inputenc}
\usepackage[margin=1in]{geometry}
\usepackage{amsmath}
\usepackage{amssymb}

\setlength{\parindent}{0pt}
\setlength{\parskip}{0.5em}

\begin{document}

\begin{center}
    {\large\textbf{STAT GU4204 --- Statistical Inference}}\\[2pt]
    {\normalsize Homework 3}\\[2pt]
    {\small Columbia University, Spring 2026 \quad $\cdot$ \quad Due: April 28, 2026 at 3pm}
\end{center}

\vspace{1em}

\section*{Section 8.5 --- Confidence Intervals}

\textbf{Exercise 8.5.2.} A random sample of 8 observations is taken from the normal distribution with unknown mean $\mu$ and unknown variance $\sigma^2$. The observed values are
\[
3.1,\ 3.5,\ 2.6,\ 3.4,\ 3.8,\ 3.0,\ 2.9,\ \text{and}\ 2.2.
\]
Find the shortest confidence interval for $\mu$ with confidence coefficients
\begin{enumerate}
    \item[(a)] $0.90$,
    \item[(b)] $0.95$, and
    \item[(c)] $0.99$.
\end{enumerate}

\vspace{0.5em}

\textbf{Exercise 8.5.4.} Suppose $X_1, \dots, X_n$ form a random sample from the normal distribution with unknown mean $\mu$ and known variance $\sigma^2$. How large a sample must be taken so that there will be a confidence interval for $\mu$ with confidence coefficient $0.95$ and length less than $0.01\sigma$?

\section*{Section 8.7 --- Unbiased Estimators}

\textbf{Exercise 8.7.11.} A drug is administered to two types of animals $A$ and $B$. The mean response of both types equals an unknown $\theta$, but the variance of type $A$ is four times the variance of type $B$. Let $X_1, \dots, X_m$ be the responses of $m$ type-$A$ animals and $Y_1, \dots, Y_n$ be the responses of $n$ type-$B$ animals. Consider the estimator
\[
\hat{\theta} = \alpha \bar{X}_m + (1 - \alpha) \bar{Y}_n.
\]
\begin{enumerate}
    \item[(a)] For what values of $\alpha$, $m$, and $n$ is $\hat{\theta}$ an unbiased estimator of $\theta$?
    \item[(b)] For fixed $m$ and $n$, what value of $\alpha$ yields an unbiased estimator with minimum variance?
\end{enumerate}

\section*{Section 8.8 --- Fisher Information and Efficiency}

\textbf{Exercise 8.8.2.} Suppose $X$ has the geometric distribution with parameter $p$. (See Sec.~5.5.) Find the Fisher information $I(p)$ in $X$.

\vspace{0.5em}

\textbf{Exercise 8.8.7.} Suppose $X_1, \dots, X_n$ form a random sample from the Bernoulli distribution with unknown parameter $p$. Show that $\bar{X}_n$ is an efficient estimator of $p$.

\vspace{0.5em}

\textbf{Exercise 8.8.17.} Let $X$ have the binomial distribution with parameters $n$ and $p$. Assume that $n$ is known. Show that the Fisher information in $X$ is
\[
I(p) = \frac{n}{p(1 - p)}.
\]

\section*{Section 9.1 --- Hypothesis Testing}

\textbf{Exercise 9.1.1.} Let $X$ have the exponential distribution with parameter $\beta$. Suppose that we wish to test the hypotheses
\[
H_0 : \beta \geq 1 \quad \text{versus} \quad H_1 : \beta < 1.
\]
Consider the test procedure $\delta$ that rejects $H_0$ if $X \geq 1$.
\begin{enumerate}
    \item[(a)] Determine the power function of the test.
\end{enumerate}

\end{document}
